\scanpage{366}
\mspage{105}

si le foncteur canonique
\[
C'\mathrel{\unsure{-|\to}}C
\]
est \inserted{également tjrs} un iso. dans $\mathrm{Cat\,hot}$, ou sinon, \unsure{du|des} moins dans $\mathrm{Hot}$.

Un intérêt particulier s'attache à $\pi_0(\unsure{\mathcal C|C})$, et de plus à sa variance pour $C\deleted{\unsure{/T}}$ \inserted{sur $T$} variable de la façon suivante. Si $\deleted{\unsure{\Sigma|S}}S$ est dans $\mathrm{Cat}$, on regarde les catégories
\[
\underline{\mathrm{Hom}}_{\mathrm{Cat}}(S,C_t)\qquad(t\in T),
\]
on trouve un foncteur
\[
t\longmapsto\underline{\mathrm{Hom}}_{\mathrm{Cat}}(S,C_t),\qquad T^\circ\mathrel{-}(\mathrm{Cat}),
\]
d'où une catégorie fibrée sur $T$, que je vais désigner par
\[
\underline{\mathscr H\mathrm{om}}_T(S_T,C).
\]

(Sauf erreur, elle s'interprète comme la catégorie fibrée sur $T$ correspondant au champ \inserted{sur $\mathrm{Top}(T)=T^\wedge$} des morphismes de $T^{\unsure{\wedge|\circ}}$/champs de $S_T$ dans $C$, où $S_T=S\times T$. (considéré comme catégorie \inserted{fibrée} sur $T$ à fibre \inserted{\unsure{cst}} \deleted{constant} $\unsure{S|D}$) dans $C$. Ici j'identifie \unsure{catégorie fibrée|catégories fibrées} sur $T$, et champs sur le topos $T^\wedge$.)

Il est immédiat que \deleted{cette} la propriété d'homotopie (p. 99) est stable par passage de $\mathcal C$ à $\underline{\mathscr H\mathrm{om}}_T(S_T,C)$. On s'intéresse à la variation de
\[
\pi_0\bigl(\underline{\mathscr H\mathrm{om}}_T(S_T,\mathcal C)\bigr)
\]

% archive sheet number: 365
